\begin{table*}[t]
  \centering
  \small
  \caption{自动 operation-stack induction 的 claim-facing 证据。该表把递归形成、hidden-label 定位消融和深度调优连到同一 profiling claim；hidden labels 只在 profiling 后评分。}
  \label{tab:induction-display}
  \begin{tabular}{p{0.18\textwidth}p{0.24\textwidth}p{0.29\textwidth}p{0.21\textwidth}}
    \toprule
    证据块 & 读者问题 & 主要数字 & 支持的结论 / 边界 \\
    \midrule
    E1 recursive formation & profiler 能否在没有用户指定 field chain 的情况下形成递归 operation stack? & 729 个 operations；15 个 induced stacks；深度直方图 2:1/3:1/4:13；session-as-evidence view 有 16 个 stacks。 & 可见边界证据足以诱导参差递归的 operation-only stacks，session 只是可选 evidence field。 不能声称自动发现所有 intent boundaries。 \\
    E2 localization ablation & induced stacks 能否作为真实 hidden-label tasks 上的可见 profiler view? & 4/6 tasks 形成 variable depth，2/6 tasks 在 material stop 停止；AP 0.2762 vs hand-configured 0.3116；work@5 0.653 vs flat 1；groups 12 vs fixed-session 285。 & induction 降低 flat inspection work 和 fixed-session fragmentation，但 AP 仍弱于 hand-configured specs。 不能把它写成 task-specific profile specs 的替代品。 \\
    E3 depth actionability & induced-stack depth 是否是实际可调的 profiling surface? & query-aware median AP 在 depth 3 最高（0.2865）；median work@5 在 depth 5 最低（0.4727）；material-split AP-best depths 覆盖 2, 3, 4, 5。 & 不同目标偏好不同递归深度，因此 depth 是 profile-configuration knob。 不能声称自动 depth selector 或 analyst-productivity 改善。 \\
    \bottomrule
  \end{tabular}
\end{table*}
